%! AP Statistics — Unit 5, Lesson 5.3 — Guided Notes (Answer Key)
\documentclass[10pt]{article}
\usepackage{apstats-article}
\usepackage{apstats-key}
\newcommand{\TallMath}[1]{$\displaystyle #1\rule[-1.4em]{0pt}{3.2em}$}

\begin{document}

\pageheader{Unit 5, Lesson 5.3}{Guided Notes --- Answer Key}
\namedateperiod

\vspace{0.1in}

\begin{objectivebox}
\textbf{Learning Targets}
\begin{enumerate}
  \item I can state the Central Limit Theorem and its conditions.
  \item I can compute $\mu_{\bar{x}}$ and $\sigma_{\bar{x}}$ for a sampling distribution of $\bar{x}$.
  \item I can determine when it is appropriate to use a normal model for $\bar{x}$.
\end{enumerate}
\end{objectivebox}

\begin{vocabbox}
\begin{description}
  \item[Central Limit Theorem] \ansline{approximately normal when $n$ is sufficiently large ($n \ge 30$).}
  \item[Standard error of $\bar{x}$] $\sigma_{\bar{x}} = \sigma/\sqrt{n}$.
  \item[Independence condition] \ansline{$n \le 10\%$ of the population.}
\end{description}
\end{vocabbox}

\begin{hookbox}
\textbf{Hook:}
\begin{itemize}
  \item Single bag: \ans{12}
  \item Mean of 36 bags: \ans{$12/\sqrt{36} = 2$}
\end{itemize}
\ansline{Averaging cancels out individual extremes. Large values pair with small values, so the mean fluctuates much less than any single observation.}
\end{hookbox}

\begin{notesbox}
\textbf{CLT Three Key Facts}

\begin{tabular}{@{} p{0.3\linewidth} p{0.6\linewidth} @{}}
Center: & \ans{$\mu$} \\[8pt]
Spread: & \ans{$\sigma/\sqrt{n}$} \\[8pt]
Shape: & \ans{normal; normal} \\
\end{tabular}
\end{notesbox}

\begin{notesbox}
\textbf{Conditions — Worked Example}

\ans{(1) Random: assumed. (2) Independence: $49 \le 10\%$ of all households (millions) --- yes. (3) Large sample: $n = 49 \ge 30$ --- yes. All conditions met.}

$\mu_{\bar{x}} = $ \ans{\$58{,}000} \hspace{1cm} $\sigma_{\bar{x}} = $ \ans{$21000/\sqrt{49} = \$3{,}000$}
\end{notesbox}

\begin{practicebox}
\textbf{Quick Check}
\begin{enumerate}
  \item
        \begin{enumerate}[label=\alph*.]
          \item \ans{$\mu_{\bar{x}} = 80$}; \ans{$\sigma_{\bar{x}} = 15/\sqrt{40} \approx 2.37$}
          \item \ansline{Yes. $n = 40 \ge 30$, so by CLT the sampling distribution of $\bar{x}$ is approximately normal even though the population is left-skewed.}
        \end{enumerate}
  \item \ansline{Not necessarily. $n = 9 < 30$, and the population is not normal (it is left-skewed). We cannot use the normal model for $\bar{x}$ here.}
\end{enumerate}
\end{practicebox}

\end{document}
